% Generated draft: requires mathematical and source review.
\documentclass[11pt]{article}
\usepackage[margin=25mm]{geometry}
\usepackage{fontspec}
\setmainfont{Latin Modern Roman}
\usepackage{amsmath,amssymb,mathtools,unicode-math}
\setmathfont{Latin Modern Math}
\usepackage{enumitem,needspace,xurl,hyperref}
\setlength{\parindent}{0pt}
\setlength{\parskip}{5pt}
\setlength{\emergencystretch}{3em}
\newcommand{\N}{\mathbb{N}}
\newcommand{\Z}{\mathbb{Z}}
\newcommand{\Q}{\mathbb{Q}}
\newcommand{\R}{\mathbb{R}}
\newcommand{\C}{\mathbb{C}}
\newcommand{\F}{\mathbb{F}}
\newcommand{\A}{\mathbb{A}}
\newcommand{\PP}{\mathbb{P}}
\newcommand{\E}{\mathbb{E}}
\newcommand{\eps}{\varepsilon}
\newcommand{\dd}{\,\mathrm{d}}
\begin{document}
\section*{UnsolvedMath ID: 11000213}
\textbf{Source problem code:} catalogue entry
\subsection{Definitions and mathematical statement}
In the source's translation/half-translation-surface setting, the Veech group is the derivative group of orientation-preserving affine automorphisms, in its chosen \(SL(2,\R)\) or projective convention. The first existence question asks for a surface with
\[\Gamma_{\mathrm{Veech}}=\langle A\rangle\cong\Z,\qquad A\text{ hyperbolic}.\]
For a projective matrix, hyperbolicity means a lift has \(|\operatorname{tr}A|>2\). The generator is hyperbolic, not parabolic, and the requested group is actually cyclic rather than merely virtually cyclic.

The source gives an equivalent formulation: find a pseudo-Anosov mapping class \(f\) whose associated Teichm\"uller disk \(D\) satisfies
\[\operatorname{Stab}_{\operatorname{Mod}(S)}(D)=\langle f\rangle.\]
The stabilizer is setwise; powers of \(f\) preserve the disk while generally acting nontrivially on it. This excludes additional stabilizing mapping classes or roots enlarging the stated cyclic group.

The asserted equivalence requires conventions. The disk stabilizer maps to its projective derivative group, and kernel automorphisms or the distinction between \(SL(2,\R)\) and \(PSL(2,\R)\) may contribute finite symmetries. A cyclic projective image alone does not automatically identify the full mapping-class stabilizer with \(\Z\). The source does not specify whether these kernels are absent or quotiented out. Both formulations are therefore retained with this compatibility gap explicit. The surface genus, compactness/finite-type, marked points and translation-versus-quadratic-differential conventions are also unstated. The summary's tentative suggestion of a parabolic-only question conflicts with the actual hyperbolic generator in the source and is not adopted.
\subsection{Short English statement}
Does some Veech group consist exactly of powers of one hyperbolic element? The source also asks for a pseudo-Anosov disk with only its cyclic stabilizer, subject to checking finite-kernel and projectivization conventions.
\subsection{Sources}
\begin{itemize}
\item Local UnsolvedMath registry, permanent ID \(11000213\), exact source code \texttt{catalogue entry}.
\item Farb, \emph{Problems on Mapping Class Groups and Related Topics}, Problem 6, ``Purely cyclic,'' PDF page 259: \url{https://www.math.uchicago.edu/~farb/papers/mcgbook.pdf}. Registry access 2026-07-29; pdf-font-structure extraction. No separate primary example or resolving reference is supplied.
\end{itemize}
\end{document}
